\chapter{总结与展望}\label{chap:conclusion}

\section{总结}
本毕业设计以电缆通道检查机器人为主题展开研究，从运动控制仿真到路径规划算法对比，完整地仿真了电缆通道检测机器人在工作时的场景。

在运动控制方面，本毕业设计在matlab的simulink仿真平台上进行机器人的整体建模与障碍物环境搭建，同时设计了lqr串联VMC的运动控制器，并验证了该运动控制器与参数配置可以满足基本运动控制与避障需求。

在路径规划方面，本毕业设计依托开源的Sparrow仿真平台进行仿真实验测试，对比了传统路径规划算法DWA与经典强化学习算法，经典强化学习算法中，DQN算法只可处理离散动作空间，遂对比测试SAC处理连续动作空间以便于在后续运动控制中机器人动作平滑不易发散。同时，SAC算法在训练轮数增长的过程中表现相较于DQN更加稳定，收敛较快且后续未出现性能震荡。

最后在SAC的基础上串联Transformer架构，融合历史雷达信息与当前状态信息一起训练，以增强模型的长时序特征提取能力。具体对应到实际中的表现为对于动态障碍物有较好的预测能力，不会像传统强化学习算法一样动态障碍物前一帧处于雷达视野而后一帧未被雷达扫描到则判定为无动态障碍物。TransSAC收敛速度明显快于SAC，在相同训练次数的情况下，TransSAC的各项指标均表现更优。

\section{展望}
在运动控制方面，本毕业设计仿真平台为matlab的simulink，后续可以考虑将仿真平台迁移到Gazebo等更贴近实际应用的仿真环境中，便于后续SimtoReal的迁移验证。同时，当前的控制器为lqr控制器，无法对输出量进行限幅，在极端工况下仍有整车发散的可能。后续可考虑MPC控制器处理带约束的问题。

在路径规划方面，本毕业设计的路径规划场景为二维地图，后续可以考虑扩展至三维地图，与实际电缆通道检测更贴近。同时，当前的路径规划算法输出线速度、角速度需进一步做最大加速度、角加速度的限制才可部署至轮腿构型的机器人，避免出现由于加速度过大整车发散的问题。后续可以考虑在状态空间中直接输出线加速度和角加速度进行训练，或者在算法输出后增加滤波模块后再输入给机器人。